\documentclass[a4paper,12pt]{article}
\usepackage[utf8]{inputenc}
\usepackage[T1]{fontenc}
\usepackage[ngerman]{babel}
\usepackage{amsmath}
\usepackage{amssymb}
\usepackage{enumitem}
\usepackage{geometry}
\geometry{a4paper, margin=2.5cm}

\title{Einsendeaufgaben -- Lektion 2\\[0.5em]
\large Modul 61112: Lineare Algebra}
\author{}
\date{}

\begin{document}
\maketitle

\section*{Aufgabe 2.2}

\begin{enumerate}[label=\alph*)]

\item
\[
  f(x,y) = \begin{pmatrix} x+y \\ (1+i)x - y \end{pmatrix}
  = x \begin{pmatrix} 1 \\ 1+i \end{pmatrix} + y \begin{pmatrix} 1 \\ -1 \end{pmatrix}
  = \begin{pmatrix} 1 & 1 \\ 1+i & -1 \end{pmatrix} \begin{pmatrix} x \\ y \end{pmatrix}
\]

\[
  \Rightarrow \; {}_{\mathcal{E}}M_{\mathcal{E}}(f) = \begin{pmatrix} 1 & 1 \\ 1+i & -1 \end{pmatrix}
\]

\item Ermittlung des Ranges von ${}_{\mathcal{E}}M_{\mathcal{E}}(f)$:

\[
  \begin{pmatrix} 1 & 1 \\ 1+i & -1 \end{pmatrix}
  \xrightarrow{Z_{21}(-1-i)}
  \begin{pmatrix} 1 & 1 \\ 0 & -2-i \end{pmatrix}
  \xrightarrow{Z_2\!\left(\frac{1}{5}(-2+i)\right)}
  \begin{pmatrix} 1 & 1 \\ 0 & 1 \end{pmatrix}
  \xrightarrow{Z_{12}(-1)}
  \begin{pmatrix} 1 & 0 \\ 0 & 1 \end{pmatrix}
\]

Die Treppennormalform von ${}_{\mathcal{E}}M_{\mathcal{E}}(f)$ ist $I$.
Folglich ist ${}_{\mathcal{E}}M_{\mathcal{E}}(f)$ invertierbar und demnach ist $f$ bijektiv.
Da $f$ auch linear ist, ist $f$ ein Isomorphismus.

\item Nach der Transformationsformel gilt:
\[
  {}_{\mathcal{B}}M_{\mathcal{B}}(f)
  = {}_{\mathcal{B}}M_{\mathcal{E}}(\mathrm{id}_{\mathbb{C}^2})
    \cdot {}_{\mathcal{E}}M_{\mathcal{E}}(f)
    \cdot {}_{\mathcal{E}}M(\mathrm{id}_{\mathbb{C}^2})_{\mathcal{B}}
\]

Nach a) gilt:
\[
  {}_{\mathcal{E}}M_{\mathcal{E}}(f) = \begin{pmatrix} 1 & 1 \\ 1+i & -1 \end{pmatrix}
\]

Nach Definition von $\mathcal{B}$ gilt:
\[
  {}_{\mathcal{E}}M_{\mathcal{B}}(\mathrm{id}_{\mathbb{C}^2}) = \begin{pmatrix} 0 & 1 \\ i & 2i \end{pmatrix}
\]

Da ${}_{\mathcal{B}}M_{\mathcal{E}}(\mathrm{id}_{\mathbb{C}^2}) = \bigl({}_{\mathcal{E}}M_{\mathcal{B}}(\mathrm{id}_{\mathbb{C}^2})\bigr)^{-1}$,
bestimmen wir die Inverse von ${}_{\mathcal{E}}M_{\mathcal{B}}(\mathrm{id}_{\mathbb{C}^2})$:

\[
  \left(\begin{array}{cc|cc} 0 & 1 & 1 & 0 \\ i & 2i & 0 & 1 \end{array}\right)
  \xrightarrow{Z_{12}}
  \left(\begin{array}{cc|cc} i & 2i & 0 & 1 \\ 0 & 1 & 1 & 0 \end{array}\right)
  \xrightarrow{Z_1(-i)}
  \left(\begin{array}{cc|cc} 1 & 2 & 0 & -i \\ 0 & 1 & 1 & 0 \end{array}\right)
  \xrightarrow{Z_{12}(-2)}
  \left(\begin{array}{cc|cc} 1 & 0 & -2 & -i \\ 0 & 1 & 1 & 0 \end{array}\right)
\]

\[
  \Rightarrow \; {}_{\mathcal{B}}M_{\mathcal{E}}(\mathrm{id}_{\mathbb{C}^2}) = \begin{pmatrix} -2 & -i \\ 1 & 0 \end{pmatrix}
\]

\[
  {}_{\mathcal{B}}M_{\mathcal{B}}(f)
  = \begin{pmatrix} -2 & -i \\ 1 & 0 \end{pmatrix}
    \begin{pmatrix} 1 & 1 \\ 1+i & -1 \end{pmatrix}
    \begin{pmatrix} 0 & 1 \\ i & 2i \end{pmatrix}
  = \begin{pmatrix} -2 & -i \\ 1 & 0 \end{pmatrix}
    \begin{pmatrix} i & 1+2i \\ -i & 1-i \end{pmatrix}
  = \begin{pmatrix} -1-2i & 1-5i \\ i & 1+2i \end{pmatrix}
\]

\item In c) haben wir schon ermittelt, dass
\[
  \begin{pmatrix} 0 & 1 \\ i & 2i \end{pmatrix}^{-1} = \begin{pmatrix} -2 & -i \\ 1 & 0 \end{pmatrix}
\]

Also ist $\mathcal{B}^* = \bigl((-2,\,-i),\,(1,\,0)\bigr)$.

\end{enumerate}

\end{document}
